%==============================================================================
% Importando a classe
%==============================================================================
\documentclass{templufla}          

%==============================================================================
% Pacote hyperref para referências cruzadas dentro do documento \ref e \autoref
% O pacote hyperref precisa ser importado o mais cedo possível.

% Setspace precisa estar antes do hyperref para usar a versão atual do abntex.
\usepackage{setspace}

% Definindo cores para os links cruzados
\usepackage{xcolor}              % para mais cores
\definecolor{rltred}{rgb}{0.2,0,0}
\definecolor{rltgreen}{rgb}{0,0.2,0}
\definecolor{rltblue}{rgb}{0,0,0.2}


% Importa hyperref
\usepackage[%
hidelinks, % colorlinks=true, | Troque se quiser link com cor
urlcolor=rltblue,       % \href{...}{...} external (URL)
filecolor=rltgreen,     % \href{...} local file
linkcolor=rltred,       % \ref{...} and \pageref{...}
citecolor=rltgreen,%
]{hyperref} % 


%==============================================================================
% Configurando a linguagem do template

% Português
\usepackage[brazilian]{babel}
\input{linguas/tempPor}

% Inglês
%\usepackage[english]{babel}
%\input{linguas/tempEng}

%==============================================================================
% Configuração dos pacotes necessários
% Primeiros pacotes (alguns precisam ser importados depois)

% TEXTOS
\usepackage[T1]{fontenc}         % usa fontes postscript com acentos
\usepackage[utf8]{inputenc}      % permite edição direta com acentos
\usepackage{listings}            % para inserção de código
\usepackage{fancyvrb}            % para inserção de saídas de comandos

% MATEMÁTICA
\usepackage{amsmath}             % pacote da AMS para Matemática Avançada
\usepackage{amssymb}             % símbolos extras da AMS
\usepackage{mathtools}			 % para mais comandos de matemática
\usepackage{latexsym}            % símbolos extras do LaTeX

% TABELAS
\usepackage{longtable}           % para tabelas muito grandes
\usepackage{array}               % para mais opções em tabelas
\usepackage{colortbl}            % para cores em tabelas
\usepackage{multirow}            % para células com várias linhas
\usepackage{multicol}            % para células com várias colunas

% FIGUAS E GRÁFICOS
\usepackage{subfigure}           % para figuras dentro de figuras
\usepackage{booktabs}
\usepackage{tikz}
\usetikzlibrary{arrows.meta, positioning, calc, fit, backgrounds}
\usepackage{pgfplots}
\pgfplotsset{compat=1.18}


%==============================================================================
% Citações/referências no formato ABNT, utilizando o pacote abntex2.
% O pacote já vem na maioria das distribuições (texLive, MikTeX, Overleaf e etc...)
% mas, caso você seja super roots e não use os listados, precisa instalar o abntex2:
% http://abntex.codigolivre.org.br/

% A versão disponível com esse template sofreu algumas alterações para se adaptar ao
% padrão atual, já que a original está desatualizada.
% Para mais informações, vizite o repositório oficial do template.


% abnt-emphasize=bf coloca o título das bibliografias em negrito
% abnt-etal-list=3 limita o número de autores antes de et al. para 3
% abnt-url-package=xurl diz para usar o pacote url
\usepackage[alf, abnt-etal-list=3, abnt-emphasize=bf]{abntex2cite}

%==============================================================================
% Ajustando os pacotes importados

% Definindo cores para elementos tabulares
\newcolumntype{Z}{|>{\columncolor[gray]{0.9}}l|} %cor cinza em células
\newcolumntype{L}[1]{>{\raggedright\let\newline\\\arraybackslash\hspace{0pt}}m{#1}}
\newcolumntype{C}[1]{>{\centering\let\newline\\\arraybackslash\hspace{0pt}}m{#1}}
\newcolumntype{R}[1]{>{\raggedleft\let\newline\\\arraybackslash\hspace{0pt}}m{#1}}

% Definindo espaçamento vertical de elementos tabulares
\renewcommand{\arraystretch}{1.5}

% Configuração padrão do listings (ambiente de código)
\lstset{
	extendedchars=true,
	breaklines=true,
	tabsize=3,
	basicstyle=\linespread{1}\ttfamily\normalsize,
	stringstyle=\em,
	showstringspaces=true,
	inputencoding=utf8,
	captionpos=t,
	backgroundcolor=\color{black!5},
	literate={{á}{{\'a}}1 {é}{{\'e}}1 {í}{{\'i}}1 {ó}{{\'o}}1 {ú}{{\'u}}1
	{Á}{{\'A}}1 {É}{{\'E}}1 {Í}{{\'I}}1 {Ó}{{\'O}}1 {Ú}{{\'U}}1
	{â}{{\^a}}1 {ê}{{\^e}}1 {ô}{{\^o}}1
	{Â}{{\^A}}1 {Ê}{{\^E}}1 {Ô}{{\^O}}1
	{ã}{{\~a}}1 {õ}{{\~o}}1 {ẽ}{{\~e}}1 
	{ç}{{\c{c}}}1
	{à}{{\`{a}}}1{À}{{\`{A}}}1}
}

% Carregamento do glossário, listas de siglas, símbolos e etc
\loadglsentries{glossarios/abreviaturas}
\loadglsentries{glossarios/glossario}
\loadglsentries{glossarios/siglas}
\loadglsentries{glossarios/simbolos}

% Especificando hifenizações que por ventura LaTeX não saiba fazer
% Por padrão 99,9% dos termos em português devem ser hifenizados corretamente.
% Adicione aqui palavras precise especificar como hifenizar (ou não hifenizar)
\hyphenation{software am-bien-te coor-de-na-ção FAE-PE TelEduc Williams UFLA Murray Bookchin}

%==============================================================================
% Dados da monografia, capa: autor, titulo, banca, etc... - SUBSTITUA DE ACORDO
%==============================================================================
% Dados para o arquivo PDF (metadados)

\hypersetup{%
% Titulo
pdftitle={SIM --- Sistema de Informações Municipais},
pdfauthor={Rafael Alves Silva Rezende},
pdfsubject={Trabalho de Conclusão de Curso -- Ciência da Computação -- UFLA},
pdfkeywords={SIM; inteligência artificial; LLM; RAG; Text-to-SQL; transparência pública; dados abertos municipais.}
}
	
%==============================================================================
% Informações básicas

% Nome da(s)/do(s) autora(s)/autor(es)
\author{Rafael Alves Silva Rezende}

\title{SIM --- Sistema de Informações Municipais:\\
Orquestração Cognitiva com RAG e Text-to-SQL Aplicada a Dados de Transparência Municipal}

\subtitle{Instância piloto em Caeté (MG)}

\engtitle{Municipal Information System (SIM): Cognitive Orchestration with RAG and Text-to-SQL}

\engsubtitle{Pilot instance for Caete, Minas Gerais, Brazil}

\date{2026}

\tipo{Trabalho de Conclusão de Curso apresentado à Universidade Federal de Lavras, como parte das exigências do Curso de Bacharelado em Ciência da Computação, para obtenção do título de Bacharel.}

\orientador{Prof. Dr. Denilson Alves Pereira}
%\orientadora{Profa. Dr. Maria Orientadora}

% Coorientador/coorientadora
% Altere o comando para os diferentes gêneros
% Comente se não houver.
%\coorientadora{Profa. Dra. Maria Orientadora }
%\coorientador{Prof. Dr. José Orientador}

% Local
\local{Lavras -- MG}

%==============================================================================
% Informações da banca

% Membros da banca. Até quatro membros. 
% Deve ser {nome}{instituição}.
% Comente os últimos se a banca tiver menos membros.
\bancaum{Prof. Dr. [Nome do Membro da Banca 1]}{[Instituição de origem]}
\bancadois{Prof. Dr. [Nome do Membro da Banca 2]}{[Instituição de origem]}
\bancatres{Prof. Dr. [Nome do Membro da Banca 3]}{[Instituição de origem]}
\bancaquatro{}{}

% Data da defesa
\defesa{[Data da defesa]}

%==============================================================================
% Ficha catalográfica, gerada pelo sistema da Biblioteca da UFLA
% https://bibliotecauniversitaria.ufla.br/recursos-tecnologicos/ficha-catalografica
% adicione o arquivo PDF da ficha ao projeto e ajuste o nome no comando

% 
%\fichaCatalograficaPDF{ficha-catalografica-SIG.pdf}

%==============================================================================
% Se vc já defendeu e tem o arquivo escaneado da folha de rosto, 
% descomente e altere o nome do arquivo
%\folhaAprovacaoAssinada{folharosto}


%==============================================================================
% O documento em si
%==============================================================================
\begin{document}

%==============================================================================
%Elementos pré-textuais

% Cria capa, folha de rosto, folha de aprovação e ficha
\maketitle

% Dedicatórias\\
\dedic{Espaço reservado à dedicatória.}

% Agradecimentos
\thanks{Espaço reservado aos agradecimentos.}         

% Epígrafe
\epigrafe{%
Espaço reservado à epígrafe.\par
(Autor Desconhecido)}

% Resumo e palavras-chave
% As palavras-chave PRECISAM vir primeiro
\palchaves{SIM; inteligência artificial; modelos de linguagem; RAG; Text-to-SQL; transparência pública; dados abertos municipais.}
\resumo{%
Este trabalho apresenta o \textbf{SIM} (\textit{Sistema de Informações Municipais}), plataforma conversacional multi-instância para consulta a dados públicos municipais no contexto da tecnologia cívica (\textit{GovTech}), com instância piloto em Caeté (MG) via WhatsApp. O sistema integra recuperação documental sobre o Diário Oficial e Text-to-SQL sobre bases SICOM/CGU (2020--2025). Descrevem-se a arquitetura de orquestração, o \textit{pipeline} de ingestão de dados e técnicas de engenharia de contexto. Para comprovar a viabilidade funcional, conduziu-se avaliação reprodutível em três lotes experimentais (o Lote I compreende quatro experimentos, E1--E4): Text-to-SQL com 80 casos executáveis derivados de dez intenções analíticas (baseline E1: VSR 97,5\% e EF 62,5\%; Prompt B com regras de domínio: VSR 100\% e EF 72,5\%) e recuperação documental com 40 perguntas estratificadas (Recall@5 de 100\% na configuração híbrida, ante 22,5\% na busca densa isolada, evidenciando a necessidade de recuperação híbrida para atos numerados; MRR 97,5\%; IC 95\% do Recall: 91,2--100\%). Artefatos reprodutíveis estão disponíveis em repositório público.
}

% Abstract e keywords
% As keywords PRECISAM vir primeiro
\keywords{Municipal Information System; artificial intelligence; large language models; RAG; Text-to-SQL; open government data.}
\abstract{%
This work presents SIM (\textit{Sistema de Informações Municipais}), a multi-instance conversational platform for municipal open government data, with a pilot deployment in Caete (Brazil) via WhatsApp. The system integrates document retrieval over the official gazette and Text-to-SQL over SICOM/CGU datasets (2020--2025). Functional viability is demonstrated through three experimental batches (Batch I comprises four experiments, E1--E4): Text-to-SQL on 80 executable cases derived from ten analytical intents (baseline E1: 97.5\% VSR and 62.5\% EF; domain-rules Prompt B: 100\% VSR and 72.5\% EF) and document retrieval on 40 stratified queries (100\% Recall@5 with hybrid retrieval versus 22.5\% with dense retrieval alone, evidencing the need for hybrid search on numbered acts; 97.5\% MRR; 95\% CI for Recall: 91.2--100\%). Reproducible artifacts are publicly available.
}

% Indicadores de impacto
\indicadoresdeimpacto{%
Os impactos descritos a seguir referem-se ao \textbf{SIM} (\textit{Sistema de Informações Municipais}), instância piloto de consulta conversacional a dados públicos de Caeté (MG) via WhatsApp. A maior parte dos efeitos é \textbf{potencial}, pois a validação quantitativa concentrou-se em ambiente controlado (\textit{benchmark} de 80 casos SQL e 40 perguntas RAG, além de amostra de 30 respostas avaliadas por juiz automático), e não em estudo de adoção em larga escala junto à população.

\textbf{Impactos sociais.} O sistema reduz a barreira técnica para que cidadãos, jornalistas e organizações da sociedade civil consultem despesas, licitações, folha de pagamento, transferências federais e atos do Diário Oficial em linguagem natural, canal já amplamente utilizado no Brasil. Contribui para o exercício do direito de acesso à informação e para o controle social sobre a gestão pública municipal. O território diretamente contemplado é o município de Caeté (código IBGE 3110004); o público beneficiado inclui moradores, servidores municipais e atores do jornalismo cívico. Os impactos sociais reportados baseiam-se no potencial de redução de barreira técnica demonstrado pela validação funcional; não foram coletadas métricas de adoção em larga escala.

\textbf{Impactos tecnológicos (concretos).} Foram entregues: (i) arquitetura integrando orquestração por LLM, RAG sobre Qdrant e Text-to-SQL sobre DuckDB; (ii) processo ETL reprodutível para bases SICOM/CGU; (iii) conjuntos de referência com 80 casos SQL e 40 perguntas RAG validadas; (iv) protocolo de avaliação com VSR, EF, Recall@$k$, MRR e JAR, organizado em três lotes experimentais (o Lote I compreende quatro experimentos, E1--E4); e (v) repositório público para replicação. Esses artefatos podem ser replicados por outras prefeituras com dados abertos semelhantes.

\textbf{Impactos econômicos (potenciais).} A solução utiliza dados públicos já disponibilizados e componentes de código aberto, reduzindo a necessidade de infraestrutura proprietária para consulta conversacional à transparência municipal. Não há, neste trabalho, análise de retorno financeiro nem contrato de prestação de serviço.

\textbf{Caráter extensionista.} O trabalho possui viés extensionista em potencial ao aplicar pesquisa em Ciência da Computação a um problema real de transparência municipal, com dados públicos já disponibilizados pelo governo local. A participação da sociedade externa ocorre indiretamente, como público-alvo do protótipo e como beneficiária de dados abertos preexistentes; não houve projeto extensionista formal registrado na UFLA nem equipe multidisciplinar além do autor e do orientador.

\textbf{Áreas temáticas da Política Nacional de Extensão.} Comunicação (divulgação de dados públicos em linguagem acessível); Direitos humanos e justiça (fomento à transparência e accountability); Tecnologia e produção (desenvolvimento de solução GovTech de código aberto).

\textbf{Alinhamento com ODS da ONU.} ODS~16 (Paz, justiça e instituições eficazes --- transparência e acesso à informação); ODS~9 (Indústria, inovação e infraestrutura --- inovação em serviços públicos digitais); ODS~11 (Cidades e comunidades sustentáveis --- governança urbana informada).
}

% Impact indicators
\impactindicators{%
The impacts below refer to SIM (\textit{Sistema de Informações Municipais}), a pilot conversational instance for public data of Caete (Brazil) via WhatsApp. Most effects are \textbf{potential}, because quantitative validation focused on a controlled benchmark (80 SQL cases and 40 RAG queries, plus a stratified sample of 30 natural-language answers judged automatically), not on large-scale adoption studies.

\textbf{Social impacts.} The system may lower technical barriers for citizens, journalists, and civil-society organizations to inspect municipal spending, procurement, payroll, federal transfers, and official-gazette acts through natural language on a widely used channel. The directly addressed territory is Caete (IBGE code 3110004). No formal usage metrics (active users or production query volume) were collected.

\textbf{Technological impacts.} Concrete deliverables include an LLM-orchestrated architecture combining RAG (Qdrant) and Text-to-SQL (DuckDB), a reproducible ETL pipeline for SICOM/CGU datasets, an 80-query golden dataset, and a documented evaluation protocol (VSR, EF, NEA, TSA, CHS, Recall@$k$, MRR, JAR), organized in three experimental batches (Batch~I comprises four comparative experiments, E1--E4).

\textbf{Economic impacts.} The solution relies on pre-existing open data and open-source components, reducing the need for proprietary infrastructure for conversational transparency tools. No formal cost--benefit analysis or service contract is reported.

\textbf{Extension character.} The work applies computing research to municipal transparency using pre-existing open data; society participates mainly as the intended audience rather than through a formal UFLA extension project.

\textbf{National Extension Policy areas.} Communication; Human rights and justice; Technology and production.

\textbf{UN SDGs.} SDG~16 (transparent institutions), SDG~9 (innovation in public digital services), SDG~11 (informed urban governance).
}

% Listas gerais
% Comente para remover.
\listoffigures  	% Lista de Figuras
%\listofilustracoes % Lista de Ilustrações
%\listofgraficos	% Lista de Gráficos
\listoftables       % Lista de Tabelas
%\listofquadros		% Lista de Quadros (sem quadros no documento)
%\listofexemplos    % Lista de Exemplos
%\listofteoremas    % Lista de Teoremas
\listofabreviaturas % Lista de Abreviaturas
\glsaddall[types={\acronymtype}]
\listofsiglas       % Lista de Siglas
\listofsimbolos     % Lista de Símbolos

% Sumário
\tableofcontents
\clearpage

%==============================================================================
% Elementos textuais

% Começa a mostrar a numeração
\pagestyle{ufla}


% Incluindo as seções
% É recomendado que as seções sejam organizadas em arquivos separados
% e incluídas no documento com o comando \include
\chapter{\MakeUppercase{Introdução}}

\section{Contextualização}

Os dados públicos municipais constituem patrimônio informacional essencial à fiscalização social, à prestação de contas e à tomada de decisão democrática. A Lei de Acesso à Informação (Lei nº~12.527/2011) consolidou o direito de acesso, mas o uso efetivo permanece restrito a usuários com conhecimento técnico em portais de transparência, formatos tabulares e linguagens de consulta estruturada.

Nesse contexto, os Modelos de Linguagem de Grande Escala (\gls{llm}) permitem interfaces conversacionais capazes de interpretar perguntas em linguagem natural e traduzi-las em operações sobre bases governamentais \cite{Yu2018Spider,Li2024BIRD}. A combinação de Recuperação Aumentada por Geração (\gls{rag}) para documentos normativos e Text-to-SQL para consultas analíticas configura uma arquitetura híbrida com potencial de ampliar o acesso a dados de transparência \cite{Lewis2020RAG}.

Este trabalho apresenta o \textbf{SIM --- Sistema de Informações Municipais}, plataforma conversacional \textbf{multi-instância} voltada à consulta de dados públicos municipais. Cada município pode operar sua própria instância com bases e persona localizadas. A \textbf{instância de referência} documentada neste trabalho é a de \textbf{Caeté (MG)} (código IBGE/SICOM 3110004), integrada ao WhatsApp por meio de fluxos automatizados e consultando dados do SICOM, da CGU (2020--2025) e do Diário Oficial municipal via busca vetorial.

\section{Problema de pesquisa}

Apesar da abertura crescente de bases governamentais --- como o SICOM do TCE-MG \cite{TCEMG_SICOM,Lei2024SICOM} --- usuários leigos enfrentam barreiras técnicas para formular consultas sobre despesas, receitas, folha de pagamento, licitações, contratos e programas sociais. Simultaneamente, benchmarks acadêmicos de Text-to-SQL \cite{Yu2018Spider} empregam métricas de avaliação que podem subestimar a utilidade de respostas em esquemas reais com aliases arbitrários e múltiplas formulações SQL semanticamente equivalentes.

\section{Objetivo geral}

Desenvolver e avaliar empiricamente o SIM, arquitetura de orquestração cognitiva baseada em \gls{llm}s integrando \gls{rag} e Text-to-SQL, com instância piloto em Caeté (MG) e protocolo de avaliação reprodutível sobre dados abertos municipais.

\section{Objetivos específicos}

\begin{itemize}
    \item Construir \textit{pipeline} de ingestão e normalização de dados SICOM/CGU para consulta analítica;
    \item Implementar subsistema Text-to-SQL com minificação de esquema e amostragem inteligente;
    \item Implementar subsistema de recuperação documental sobre o Diário Oficial;
    \item Orquestrar os fluxos em ambiente de produção via mensageria instantânea;
    \item Definir conjunto de referência (\textit{golden dataset}) e conduzir avaliação formal com métricas de execução v3, validade SQL, latência e custo;
    \item Realizar estudos de ablação comparando variantes de \textit{prompt} e modelos de linguagem.
\end{itemize}

\section{Justificativa}

A solução proposta alinha-se à governança digital e ao ODS~16, ao reduzir a barreira técnica de acesso à informação pública. A avaliação quantitativa em dados municipais reais --- com métricas adaptadas ao contexto de perguntas abertas --- contribui com evidências empíricas para decisões de arquitetura em soluções de tecnologia cívica (\textit{GovTech}).

\section{Organização do trabalho}

O Capítulo~2 apresenta o referencial teórico. O Capítulo~3 descreve metodologia e arquitetura. O Capítulo~4 analisa resultados. O Capítulo~5 discute limitações e implicações. O Capítulo~6 encerra com conclusões.

\chapter{\MakeUppercase{Referencial Teórico}}

\section{Modelos de linguagem e dados governamentais}

Os Modelos de Linguagem de Grande Escala (\gls{llm}) generalizam padrões linguísticos e realizam raciocínio aproximado sobre instruções em linguagem natural \cite{OpenAI2024GPT4}. Em bases governamentais heterogêneas, o desafio central é condicionar o modelo com contexto suficiente --- esquema relacional, amostras de valores e regras de domínio --- sem exceder limites de tokens e custo operacional \cite{Liu2023LostMiddle}.

Diferentemente de consultas ad hoc em portais de transparência, interfaces conversacionais exigem interpretação de intenção, desambiguação de termos fiscais e seleção dinâmica entre fontes tabulares e documentais. A literatura de interfaces em linguagem natural para bases de conhecimento \cite{KatsogiannisMeimaris2023GovChat} e de \textit{GovTech} \cite{Janssen2020GovTech} aponta que a barreira de acesso não é apenas técnica, mas também semântica: o cidadão formula perguntas em prosa, enquanto os dados públicos estão organizados em esquemas normalizados e vocabulários controlados.

\section{Recuperação aumentada por geração}

\subsection{Pipeline RAG}

A técnica \gls{rag}, formalizada por \citeonline{Lewis2020RAG}, complementa o \gls{llm} com recuperação de trechos documentais relevantes antes da geração da resposta. O \textit{pipeline} clássico compreende: (i)~indexação de documentos em unidades recuperáveis (\textit{chunks}); (ii)~vetorização por modelo de \textit{embedding}; (iii)~busca por similaridade no índice; e (iv)~condicionamento do gerador com os trechos recuperados, conforme o survey de Gao, Y. et al. \cite{Gao2024RAGSurvey}.

No SIM, o mecanismo é adequado a perguntas sobre leis, decretos e normas municipais indexadas a partir do Diário Oficial. A qualidade da recuperação depende criticamente do \textit{chunking} --- tamanho e fronteira dos trechos --- e da deduplicação de conteúdo repetido no índice, que pode inflar o top-$k$ com variantes do mesmo parágrafo \cite{Asai2024SelfRAG}.

\subsection{Embeddings densos e busca esparsa}

A recuperação densa projeta consulta e documentos em espaço vetorial contínuo, capturando similaridade semântica \cite{Karpukhin2020DPR}. Modelos como \texttt{text-embedding-3-small} oferecem bom custo-benefício para índices de médio porte. Entretanto, consultas com identificadores explícitos --- número e ano de ato normativo --- frequentemente exigem correspondência lexical exata que embeddings não preservam de forma confiável \cite{Robertson2009BM25}.

A busca esparsa (\gls{bm25}) modela relevância lexical e permanece competitiva em cenários com termos raros ou identificadores numéricos \cite{Robertson2009BM25}. A \textbf{busca híbrida} combina escores densos e esparsos --- por exemplo, fusão por soma ponderada ou \textit{reciprocal rank fusion} --- para aproveitar complementaridade entre semântica e léxico \cite{Lin2021Hybrid,Lin2021Pyserini}.

\subsection{Reranking, filtros e deduplicação}

Após a recuperação inicial, um estágio de \textit{reranking} pode reordenar candidatos com modelo cruzado consulta-documento, elevando precisão no top-$k$ \cite{Nogueira2019MonoT5}. Em domínios normativos, metadados estruturados --- tipo de ato, número, ano --- permitem filtros determinísticos que reduzem falsos positivos antes da geração \cite{Gao2024RAGSurvey}.

No contexto do SIM, a deduplicação de \textit{page\_content} idêntico e o filtro por \texttt{tipo\_ato} atuam como pós-processamento leve, sem reranker neural adicional, mantendo latência compatível com canal conversacional síncrono.

\section{Text-to-SQL}

\subsection{Formulação do problema}

Text-to-SQL converte perguntas em consultas estruturadas sobre bancos relacionais. O desafio envolve \textit{schema linking} (mapear entidades da pergunta a tabelas e colunas), composição de \textit{joins}, agregações e filtros compatíveis com o dialeto SQL do \gls{sgbd} \cite{Yu2018Spider,Wang2020RATSQL}.

\subsection{Métricas de avaliação: EX e variantes}

O benchmark Spider \cite{Yu2018Spider} estabeleceu o \textit{Execution Accuracy} (EX) como comparação de resultados executados entre SQL gerada e SQL de referência. \citeonline{Zhong2020TestSuite} aprimoraram a metodologia com \textit{test suites} para reduzir falsos positivos causados por bancos com múltiplas soluções válidas.

O benchmark BIRD \cite{Li2024BIRD} aproxima aplicações reais ao incorporar bases volumosas, valores ruidosos e conhecimento externo; modelos de ponta ainda distam da performance humana em EX estrita. Em Spider, sistemas estado da arte reportam EX acima de 80\% em configurações controladas; em BIRD, EX de modelos comerciais frequentemente permanece abaixo de 60\% \cite{Li2024BIRD,Pourreza2024DIN}.

Este trabalho adota a \textbf{equivalência funcional} (EF) como métrica principal, relaxando EX estrita para colunas de resposta e tolerância numérica em grandezas fiscais --- decisão justificada porque EX estrita no baseline E1 foi de apenas 10\%, enquanto a EF foi de 62,5\% (Capítulo~4).

\subsection{Métodos com LLM}

Abordagens recentes empregam decomposição de consultas, \textit{self-correction} e seleção dinâmica de exemplos \cite{Pourreza2024DIN,Gao2024DAILSQL}. DIN-SQL decompõe a geração em etapas (esquema, colunas, SQL final) com correção iterativa; DAIL-SQL, de Gao, D. et al., enfatiza \textit{schema linking} e exemplos em contexto \cite{Gao2024DAILSQL}. RAT-SQL introduz codificação relacional do esquema para melhorar ligação entre pergunta e tabelas \cite{Wang2020RATSQL}.

No SIM, a engenharia de contexto --- minificação de DDL, amostragem de valores e regras de domínio SICOM --- substitui parte da complexidade arquitetural desses métodos, adequando-se a um único município com esquema fixo e canal de baixa latência.

\section{Orquestração cognitiva}

Designa a camada que interpreta a intenção do usuário e seleciona dinamicamente o mecanismo mais adequado --- recuperação documental ou consulta tabular. Plataformas de automação (\textit{low-code}) permitem integrar APIs, modelos e canais conversacionais em arquitetura modular \cite{Russell2020AIMA}.

No SIM, o orquestrador recebe a pergunta via WhatsApp, decide entre subsistema RAG (perguntas normativas) e Text-to-SQL (perguntas analíticas sobre SICOM/CGU), e formula a resposta em linguagem natural. Essa separação evita forçar SQL para perguntas sobre decretos e evita RAG para agregações fiscais.

\section{Dados públicos municipais no Brasil}

O SICOM, mantido pelo TCE-MG, centraliza dados contábeis e de compras públicas dos municípios mineiros \cite{TCEMG_SICOM,Lei2024SICOM}. Complementarmente, a CGU e o Siconfi \cite{STN_SiconfiAPI,CGU_Transparencia} disponibilizam dados federais cruzáveis por documento de credores e programas sociais. A política de dados abertos no Brasil \cite{Brasil2016DadosAbertos} e a Lei de Acesso à Informação \cite{Lei2012LAI} fundamentam o direito de consulta, mas não garantem interfaces acessíveis ao cidadão leigo.

Essas fontes apresentam heterogeneidade de esquema, lacunas de publicação, campos em formato YYYYMMDD e volumes que desafiam tanto o processamento quanto a avaliação de sistemas conversacionais --- características compartilhadas por centenas de municípios mineiros no mesmo padrão SICOM.

\section{Trabalhos relacionados}

\subsection{Benchmarks Text-to-SQL}

Spider \cite{Yu2018Spider} e BIRD \cite{Li2024BIRD} definem protocolos com conjuntos de referência, comparação de resultados executados e divisão treino/teste por domínio. Ambos assumem esquemas em inglês e bancos sintéticos ou generalistas; o SIM opera sobre esquema SICOM real, em português, com regras contábeis locais.

DIN-SQL \cite{Pourreza2024DIN} e DAIL-SQL de Gao, D. et al. \cite{Gao2024DAILSQL} demonstram ganhos com decomposição e \textit{schema linking}, mas em \textit{benchmarks} acadêmicos com orçamento de tokens elevado. O SIM posiciona-se como sistema aplicado de GovTech municipal: prioriza latência, canal WhatsApp e regras de domínio portáveis ao ecossistema \gls{sicom}, em detrimento de arquiteturas multiagente pesadas.

\subsection{RAG em documentos governamentais}

Trabalhos de recuperação em corpora jurídicos e administrativos enfatizam metadados de citação e estrutura hierárquica de normas \cite{Gao2024RAGSurvey}. Poucos relatam avaliação sobre Diários Oficiais municipais brasileiros com índice vetorial de milhares de \textit{chunks}. O SIM contribui com protocolo estratificado (\texttt{com\_numero} \textit{versus} temático) e evidência de que busca híbrida é necessária para consultas com número explícito de ato.

\subsection{GovTech e transparência}

Experiências com dados do SICOM para análise fiscal e detecção de irregularidades \cite{Lei2024SICOM} demonstram o potencial de ciência de dados sobre bases abertas municipais. O SIM difere ao focar interface conversacional para o cidadão, não painel analítico para especialistas. A avaliação complementar com LLM-as-Judge \cite{Zheng2023LLMJudge} mede adequação da resposta em linguagem natural, camada não capturada exclusivamente por comparação sobre SQL.

\subsection{Posicionamento do SIM}

Em síntese, o SIM integra contribuições de três linhas --- Text-to-SQL com EF sobre SICOM, RAG híbrido sobre Diário Oficial e orquestração em canal conversacional --- em arquitetura multi-instância replicável. A avaliação métrica comprova viabilidade funcional na instância Caeté; não constitui, por si só, o objeto central da pesquisa, que permanece o acesso democrático a dados públicos municipais.

\chapter{\MakeUppercase{Metodologia}}

\section{Tipo de pesquisa}

Trata-se de pesquisa aplicada, de natureza experimental e desenvolvimento tecnológico, com validação funcional em ambiente real (WhatsApp) e comprovação quantitativa reprodutível do subsistema Text-to-SQL em ambiente controlado, versionado no repositório \texttt{avaliacao/}.

\section{O SIM e a instância Caeté}

O \textbf{SIM} (\textit{Sistema de Informações Municipais}) é concebido como plataforma multi-instância: a arquitetura é replicável; variam os dados, o código IBGE e eventuais lacunas de transparência local. A instância avaliada neste trabalho é Caeté (MG), código IBGE/SICOM 3110004.

As fontes incluem:

\begin{itemize}
    \item \textbf{SICOM}: receitas, despesas, empenhos, folha, frota, licitações e contratos (2020--2025);
    \item \textbf{CGU}: transferências federais e programas sociais cruzáveis por documento do credor;
    \item \textbf{Diário Oficial}: documentos normativos indexados em Qdrant para recuperação aumentada.
\end{itemize}

Um processo ETL em Python consolida os arquivos tabulares, gera metadados de esquema (\gls{ddl}) e expõe uma API analítica com DuckDB \textit{in-memory}. A prefeitura não publicou dados de folha de pagamento e cadastro de servidores para 2020--2022; essa lacuna constitui limitação documentada da base local, incorporada ao escopo de resposta do sistema.

\section{Arquitetura do sistema}

A Figura~\ref{fig:arquitetura} resume o fluxo conversacional da instância Caeté:

% Diagrama da arquitetura do SIM — incluído via \input em metodologia.tex
\begin{figure}[htbp]
    \centering
    \resizebox{\textwidth}{!}{%
    \begin{tikzpicture}[
        font=\small,
        node distance=0.55cm and 0.9cm,
        box/.style={draw, rounded corners=2pt, align=center, minimum width=2.6cm, minimum height=0.85cm, fill=black!4},
        databox/.style={box, fill=black!8, minimum width=2.4cm},
        arrow/.style={-{Stealth[length=2.2mm]}, thick},
        lbl/.style={font=\footnotesize, fill=white, inner sep=1pt}
    ]
        \node[box, minimum width=3.2cm] (user) {Usuário\\\textit{WhatsApp}};

        \node[box, right=of user] (waha) {Gateway\\WAHA};
        \node[box, right=of waha] (n8n) {Automação\\n8n};
        \node[box, right=of n8n, minimum width=3.0cm] (orch) {Agente orquestrador\\GPT-4o};

        \draw[arrow] (user) -- (waha);
        \draw[arrow] (waha) -- (n8n);
        \draw[arrow] (n8n) -- (orch);

        \node[box, below left=1.0cm and 0.2cm of orch, minimum width=3.1cm] (rag) {RAG documental\\Qdrant + Diário Oficial};
        \node[box, below right=1.0cm and 0.2cm of orch, minimum width=3.1cm] (sql) {Text-to-SQL\\GPT-4.1-mini};

        \draw[arrow] (orch.south) -- ++(0,-0.35) -| (rag.north) node[pos=0.22, lbl, above] {normativo};
        \draw[arrow] (orch.south) -- ++(0,-0.35) -| (sql.north) node[pos=0.22, lbl, above] {tabular};

        \node[databox, below=of rag] (qdrant) {Índice vetorial\\Qdrant};
        \node[databox, below=of sql] (duck) {API analítica\\DuckDB};
        \node[databox, below=1.1cm of duck, minimum width=6.8cm] (etl) {ETL --- bases SICOM, CGU e metadados de esquema};

        \draw[arrow] (rag) -- (qdrant);
        \draw[arrow] (sql) -- (duck);
        \draw[arrow] (duck) -- (etl);

        \draw[arrow] (orch.east) -- ++(0.55,0) |- (n8n.east) node[pos=0.78, lbl, right] {resposta em linguagem natural};
        \draw[arrow] (n8n.west) -- (waha.east);
        \draw[arrow] (waha.west) -- (user.east);

        \node[font=\footnotesize\itshape, below=0.35cm of etl] (leg) {Avaliação Text-to-SQL conduzida em pipeline offline (\texttt{avaliacao/}), desacoplada do canal WhatsApp.};
    \end{tikzpicture}%
    }
    \caption{Arquitetura do SIM (instância Caeté): fluxo conversacional em produção e subsistemas de recuperação documental e consulta tabular.}
    \label{fig:arquitetura}
\end{figure}


\begin{enumerate}
    \item O usuário envia a pergunta via WhatsApp;
    \item A plataforma de automação recebe o evento e aciona o agente orquestrador;
    \item O agente seleciona recuperação documental (Qdrant) ou tradução Text-to-SQL;
    \item O subsistema SQL utiliza modelo de linguagem com \textit{prompt} contendo esquema compactado e amostras de dados;
    \item O resultado tabular é interpretado pelo orquestrador e devolvido em linguagem natural ao usuário.
\end{enumerate}

\section{Engenharia de contexto}

\subsection{Minificação de esquema (\textit{DDL minification})}

Consiste na remoção de redundâncias, padronização de descrições e simplificação do esquema relacional enviado ao modelo, reduzindo tokens sem eliminar restrições semânticas críticas.

\subsection{Amostragem inteligente (\textit{smart sampling})}

Inclusão de exemplos reais de valores por coluna relevante, visando melhorar inferência de filtros, \textit{joins} e convenções locais.

\subsection{Variantes de \textit{prompt} avaliadas}

\begin{itemize}
    \item \textbf{Prompt compacto} ($\approx$16\,500 tokens): esquema minificado e amostras --- referência do Lote I;
    \item \textbf{Prompt estendido} ($\approx$35\,000 tokens): esquema com descrições ampliadas;
    \item \textbf{Prompt v2.1} ($\approx$16\,500 tokens): variante final com regras de domínio --- configuração adotada após o Lote II.
\end{itemize}

\section{Conjunto de referência (\textit{golden dataset})}

Foi construído um conjunto de 80 consultas em linguagem natural, estratificadas por dificuldade (fácil, médio, difícil) e por domínio temático. Cada item contém a pergunta, a consulta SQL de referência validada no DuckDB, metadados auxiliares e a coluna \texttt{colunas\_resposta} indicando quais grandezas respondem à pergunta (75 itens por procedimento automatizado, 5 por curadoria manual).

Dez consultas-base foram expandidas com variações paramétricas (anos, limites e agregações), totalizando 80 casos executáveis. Os artefatos estão em \texttt{avaliacao/dados/golden/}.

\section{Protocolo de validação do Text-to-SQL}

A avaliação do subsistema Text-to-SQL foi conduzida por \textit{pipeline} automatizado desacoplado do canal WhatsApp, garantindo reprodutibilidade (detalhes em Apêndice~\ref{ap:artefatos}). A validação do subsistema RAG segue protocolo complementar descrito na Seção seguinte.

\section{Protocolo de validação do RAG}

Complementarmente ao Text-to-SQL, conduziu-se avaliação da recuperação documental sobre o Diário Oficial municipal indexado em Qdrant (9\,346 \textit{chunks}, embedding \texttt{text-embedding-3-small}).

\subsection{Conjunto de referência RAG}

O \textbf{conjunto de referência RAG ampliado} contém $N_{\text{RAG}} = 40$ perguntas estratificadas em dois estratos:

\begin{itemize}
    \item \textbf{com\_numero} ($n = 32$): decretos sorteados uniformemente do inventário do índice (seed 42), com pergunta-template contendo número e ano explícitos do ato;
    \item \textbf{tematico} ($n = 8$): perguntas curadas sobre conteúdo normativo, reaproveitadas da amostra temática inicial, sem depender exclusivamente de número na formulação.
\end{itemize}

A inclusão dos itens \textit{com\_numero} é independente do \textit{rank} de recuperação --- critério de amostra honesta. Cada item possui documento de referência validado no índice (cabeçalho do ato presente em \textit{page\_content}).

\subsection{Indicadores de recuperação}

Para cada pergunta $i$, registra-se o ranking do documento esperado entre os $k$ primeiros \textit{chunks} retornados ($k = 5$):

\begin{equation}
\text{Recall@}k = \frac{1}{N_{\text{RAG}}} \sum_{i=1}^{N_{\text{RAG}}} \mathbb{1}\!\left[\text{doc esperado} \in \text{top-}k_i\right]
\end{equation}

A \textbf{taxa de recuperação média} (MRR --- \textit{Mean Reciprocal Rank}) mede a posição média do primeiro acerto:

\begin{equation}
\text{MRR} = \frac{1}{N_{\text{RAG}}} \sum_{i=1}^{N_{\text{RAG}}} \frac{1}{\text{rank}_i}
\end{equation}

onde $\text{rank}_i$ é a posição do \textit{chunk} que contém o ato esperado, ou zero se ausente do top-$k$. Reportam-se também intervalos de confiança de Wilson (95\%) \cite{Agresti1998} para as proporções de acerto.

\subsection{Escopo e limitações do protocolo RAG}

A avaliação compara estratégias de recuperação (busca densa, deduplicação, filtro por tipo de ato e busca híbrida por número/ano) sobre o mesmo \textit{golden}. Não avalia geração da resposta final em linguagem natural nem o roteamento do orquestrador.


\subsection{Notação e pré-condições}

Para cada consulta $i$ no conjunto de referência ($N = 80$):

\begin{itemize}
    \item $q_i$ --- pergunta em linguagem natural;
    \item $\text{SQL}_i^{(r)}$ --- consulta de referência (gabarito);
    \item $\text{SQL}_i^{(g)}$ --- consulta gerada pelo modelo;
    \item $R_i$ --- resultado tabular de $\text{SQL}_i^{(r)}$;
    \item $G_i$ --- resultado tabular de $\text{SQL}_i^{(g)}$;
    \item $C_i$ --- conjunto de colunas de resposta declaradas no gabarito.
\end{itemize}

Define-se a função de normalização de célula $\nu(\cdot)$, aplicada antes de qualquer comparação: nulos unificados; números arredondados em quatro casas decimais; strings sem espaços laterais e em maiúsculas; strings numéricas convertidas com vírgula substituída por ponto.

\subsection{Taxa de SQL válida (VSR)}

A \textbf{taxa de SQL válida} (\textit{Valid SQL Rate}, VSR) mede se a consulta gerada é sintaticamente correta e executável sobre o esquema municipal:

\begin{equation}
\text{VSR} = \frac{1}{N} \sum_{i=1}^{N} \mathbb{1}\!\left[\text{SQL}_i^{(g)} \text{ executa sem erro}\right]
\end{equation}

onde $\mathbb{1}[\cdot]$ é a função indicadora. Quando a SQL gerada falha, os indicadores de equivalência de resultado para a consulta $i$ são nulos.

\subsection{Equivalência funcional (EF)}

A \textbf{equivalência funcional} (EF) verifica se o cidadão receberia as grandezas corretas, independentemente de aliases de coluna ou colunas auxiliares no resultado gerado. Trata-se da métrica adotada neste trabalho para comprovar a utilidade do resultado tabular.

Diferentemente do \textit{Execution Accuracy} (EX) usado em Spider e BIRD --- que exige correspondência estrita de nomes de colunas e do multiconjunto de pares (coluna, valor) ---, a EF restringe a comparação às colunas de resposta declaradas no conjunto de referência, admite mapeamento entre aliases de coluna e aplica tolerância numérica $\tau$ em grandezas fiscais. Essa definição penaliza menos formulações SQL semanticamente equivalentes, alinhando-se ao objetivo de comprovar utilidade para o cidadão.

Para cada consulta $i$ com VSR positivo, exige-se $|R_i| = |G_i|$ (mesma cardinalidade de linhas). Para cada coluna $k \in C_i$, deve existir alguma coluna $j$ em $G_i$ cujo vetor de valores normalizados ao longo das linhas seja equivalente ao vetor da coluna correspondente em $R_i$, com tolerância numérica $\tau = 0{,}5\%$ em valores fiscais:

\begin{equation}
\left|\nu(a) - \nu(b)\right| / \max\!\left(|\nu(b)|, \varepsilon\right) \leq \tau, \quad \tau = 0{,}005
\end{equation}

onde $\varepsilon$ evita divisão por zero. A equivalência funcional agregada é:

\begin{equation}
\text{EF} = \frac{1}{N} \sum_{i=1}^{N} \mathbb{1}\!\left[\text{EF}_i = 1\right]
\end{equation}

com $\text{EF}_i = 1$ quando todas as colunas em $C_i$ possuem correspondente equivalente em $G_i$ sob as regras acima. Proporções agregadas são reportadas com intervalos de confiança de Wilson (95\%) \cite{Agresti1998}; comparações pareadas E1 \textit{versus} Prompt B utilizam teste de McNemar quando aplicável.

\subsection{Indicadores complementares}

\begin{description}
    \item[NEA, TSA, CHS] Indicadores auxiliares de resultado não vazio, seleção de tabelas e presença de filtros relevantes.
    \item[JAR --- \textit{Judge Agreement Rate}] Proporção de respostas em linguagem natural classificadas como corretas ou parcialmente corretas por LLM-as-Judge \cite{Zheng2023LLMJudge}, em amostra estratificada ($M = 30$):
    \begin{equation}
    \text{JAR} = \frac{1}{M} \sum_{j=1}^{M} \mathbb{1}\!\left[\text{veredito}_j \in \{\text{correto}, \text{parcialmente\_correto}\}\right]
    \end{equation}
\end{description}

A implementação completa, incluindo variantes diagnósticas, está documentada no repositório público de artefatos do trabalho (Apêndice~\ref{ap:artefatos}).

\subsection{Organização dos lotes experimentais}

A validação empírica organizou-se em três lotes sobre dados fixos (conjunto SQL de 80 casos; conjunto RAG ampliado de 40 perguntas):

\begin{enumerate}
    \item \textbf{Lote I} --- comparação de modelo e tamanho de contexto (E1--E4), fixando GPT-4.1-mini e Prompt A como referência;
    \item \textbf{Lote II} --- ablation de engenharia de \textit{prompt} (Prompt B), mantendo modelo e conjunto de referência do Lote I;
    \item \textbf{Lote III} --- comparação de estratégias de recuperação documental (busca densa, deduplicação, filtro por tipo de ato e busca híbrida), mantendo índice e embedding fixos.
\end{enumerate}

Cada lote isolou uma variável experimental.

\subsection{Experimentos do Lote I}

\begin{enumerate}
    \item \textbf{E1 --- Baseline}: GPT-4.1-mini com Prompt A (compacto);
    \item \textbf{E2 --- Ablação de contexto}: GPT-4.1-mini com \textit{prompt} estendido;
    \item \textbf{E3 --- Ablação de modelo (econômico)}: GPT-4o-mini com Prompt A;
    \item \textbf{E4 --- Ablação de modelo (avançado)}: GPT-4o com Prompt A.
\end{enumerate}

\section{Ambiente experimental}

\begin{itemize}
    \item DuckDB local sobre bases tabulares SICOM/CGU consolidadas;
    \item Modelos OpenAI via API oficial, temperatura 0,0;
    \item Saída estruturada em JSON contendo a consulta SQL;
    \item Persistência incremental em arquivo de \textit{checkpoint} para retomada de experimentos.
\end{itemize}


\chapter{\MakeUppercase{Análise de Resultados}}

\section{Validação do subsistema Text-to-SQL}

A Tabela~\ref{tab:v3} consolida os quatro experimentos. Os indicadores comprovam que o pipeline gera SQL executável em grande parte dos casos (VSR entre 90\% e 97,5\%) e produz resultados funcionalmente equivalentes ao gabarito na maioria das consultas do baseline.

\begin{table}[htbp]
\centering
\caption{Indicadores de validação --- comparativo dos experimentos ($n=80$)}
\label{tab:v3}
\begin{tabular}{lrr}
\toprule
\textbf{Experimento} & \textbf{VSR (\%)} & \textbf{EF (\%)} \\
\midrule
E1 --- Baseline (GPT-4.1-mini, compacto) & 97,50 & \textbf{62,50} \\
E2 --- Contexto estendido (GPT-4.1-mini) & 96,25 & 58,75 \\
E3 --- GPT-4o-mini & 90,00 & 46,25 \\
E4 --- GPT-4o & 97,50 & 57,50 \\
\bottomrule
\end{tabular}
\end{table}

\section{Configuração baseline (E1)}

O experimento E1 avaliou 80 consultas com GPT-4.1-mini e \textit{prompt} compacto --- configuração adotada em produção na instância Caeté. A Tabela~\ref{tab:baseline} resume os indicadores e a latência observada.

\begin{table}[htbp]
\centering
\caption{Resultados do experimento E1 --- baseline}
\label{tab:baseline}
\begin{tabular}{lr}
\toprule
\textbf{Indicador} & \textbf{Valor} \\
\midrule
VSR & 97,50\% (78/80) \\
Equivalência funcional (EF) & 62,50\% (50/80) \\
Latência média de geração & 6,2 s \\
Latência média de execução SQL & 1,1 s \\
\bottomrule
\end{tabular}
\end{table}

\section{Resultados por dificuldade (E1)}

\begin{table}[htbp]
\centering
\caption{E1 --- indicadores por nível de dificuldade}
\label{tab:diff}
\begin{tabular}{lrrr}
\toprule
\textbf{Dificuldade} & \textbf{$n$} & \textbf{EF (\%)} & \textbf{VSR (\%)} \\
\midrule
Fácil & 22 & 68,18 & 100,00 \\
Médio & 44 & 61,36 & 97,73 \\
Difícil & 14 & 57,14 & 92,86 \\
\bottomrule
\end{tabular}
\end{table}

\section{Comparativo de configurações}

O \textit{prompt} estendido (E2) não superou o compacto em equivalência funcional (58,75\% \textit{versus} 62,50\%). O GPT-4o-mini (E3) degrada VSR (90\%) e EF (46,25\%). O GPT-4o (E4) não supera o baseline em EF (57,50\%) e apresenta latência média de 29,5~s na geração de SQL --- inviável para respostas síncronas em canal conversacional.

\section{Refinamento do \textit{prompt} Text-to-SQL}

Após os experimentos comparativos, incorporaram-se ao \textit{prompt} regras derivadas de falhas reais do baseline e de perguntas de testers: escopo de filtros de higiene por formato de resposta, datas em formato YYYYMMDD, tabelas setoriais (SAAE, RSP), busca textual com \texttt{ILIKE} e cruzamentos por chave única. Manteve-se o modelo GPT-4.1-mini (temperatura 0) e o golden v1 ($n=80$).

\begin{table}[htbp]
\centering
\caption{Ablation de \textit{prompt} --- pós-validação ($n=80$, GPT-4.1-mini)}
\label{tab:prompt_v2}
\begin{tabular}{lrr}
\toprule
\textbf{Variante} & \textbf{VSR (\%)} & \textbf{EF (\%)} \\
\midrule
E1 --- \textit{prompt} compacto v1 (baseline) & 97,50 & 62,50 \\
e5 --- v2.0 (insights) & 95,00 & 71,25 \\
e6 --- v2.1 (insights) & \textbf{100,00} & \textbf{72,50} \\
e7 --- v2.1.1 (ajuste de escopo) & 98,75 & 72,50 \\
e8 --- v2.1 + DDL do schema real & 98,75 & 70,00 \\
\bottomrule
\end{tabular}
\end{table}

A variante v2.1 (e6) elevou a equivalência funcional em 10 pontos percentuais em relação ao E1, com VSR integral. A iteração v2.1.1 manteve EF mas reduziu VSR; o DDL extraído do schema real não superou v2.1 no golden v1 --- o gabarito foi construído sobre as convenções do esquema minificado original.

\section{Validação operacional}

Além do \textit{benchmark} offline, o SIM foi validado em ambiente operacional via WhatsApp na instância Caeté, confirmando viabilidade da orquestração integrada entre automação, mensageria, recuperação documental e consulta analítica.

\section{Avaliação da resposta em linguagem natural}

\subsection{Protocolo}

Para amostra estratificada do E1 ($n=30$), registrou-se a cadeia completa: SQL gerada, resultado tabular, resposta em linguagem natural do orquestrador (GPT-4o) e julgamento automático (GPT-4.1-mini) conforme rubrica em quatro dimensões Likert \cite{Zheng2023LLMJudge}.

\subsection{Resultados}

\begin{table}[htbp]
\centering
\caption{Avaliação LLM-as-Judge (E1, $n=30$)}
\label{tab:jar}
\begin{tabular}{lr}
\toprule
\textbf{Indicador} & \textbf{Valor} \\
\midrule
JAR & 100,00\% (30/30) \\
Adequação semântica (média) & 4,87 / 5 \\
Fidelidade aos dados (média) & 4,87 / 5 \\
Completude (média) & 4,73 / 5 \\
Clareza ao cidadão (média) & 4,77 / 5 \\
\bottomrule
\end{tabular}
\end{table}

O julgamento automático indica que as respostas em linguagem natural são adequadas ao cidadão na amostra avaliada. Essa evidência complementa, mas não substitui, a equivalência funcional sobre resultados tabulares.

\section{Validação do subsistema RAG}

Complementarmente, avaliou-se a recuperação documental sobre o Diário Oficial com o \textit{golden} RAG v1.1 ($N_{\text{RAG}} = 8$ decretos municipais com \textit{chunk} confirmado no índice Qdrant). O \textit{pipeline} gera embeddings (\texttt{text-embedding-3-small}), consulta a \textit{collection} \texttt{jornais\_caete} e compara os top-5 \textit{chunks} ao documento de referência.

\begin{table}[htbp]
\centering
\caption{Recuperação documental --- comparativo no golden v1.1 ($k=5$)}
\label{tab:rag}
\begin{tabular}{lrr}
\toprule
\textbf{Configuração} & \textbf{Recall@5 (\%)} & \textbf{MRR (\%)} \\
\midrule
Busca densa (baseline) & 100,00 (8/8) & 63,33 \\
Deduplicação de \textit{chunks} idênticos & 100,00 & 69,79 \\
Filtro por tipo de ato + dedup & 100,00 & 72,92 \\
Híbrido por nº de ato + dedup + filtro tipo & 100,00 & \textbf{87,50} \\
\bottomrule
\end{tabular}
\end{table}

O Recall@5 integral na busca densa indica que, para a amostra curada, o ato esperado aparece entre os cinco primeiros \textit{chunks}. O MRR de 63,33\% reflete que o acerto nem sempre ocorre na primeira posição --- decretos tematicamente similares competem no ranking. A deduplicação de \textit{page\_content} idêntico no top-5 elevou o MRR em 6,46 pontos percentuais sem perda de Recall. A busca híbrida antecede \textit{chunks} cujo cabeçalho cita o número e o ano do ato extraídos da pergunta, elevando o MRR a 87,50\%.

Complementarmente, avaliou-se um \textit{golden} v1.2 com $N_{\text{RAG}} = 20$ atos sorteados sem critério de inclusão por rank (amostra honesta). Na busca densa isolada, Recall@5 foi 0\%; com stack híbrido+dedup, Recall@5 e MRR atingiram 100\% --- evidência de que perguntas com número explícito de ato exigem recuperação híbrida. A amostra v1.1 permanece a referência principal do TCC; generalização e avaliação da resposta em linguagem natural permanecem como trabalhos futuros.

\chapter{\MakeUppercase{Discussão}}

\section{Acesso a dados e transparência pública}

O SIM materializa o princípio de que dados abertos só cumprem sua função democrática quando acessíveis \cite{Brasil2016DadosAbertos,Lei2012LAI}. Ao integrar consulta analítica e recuperação documental em canal conversacional amplamente utilizado no Brasil, a plataforma reduz a distância entre a publicação formal de bases governamentais e o uso efetivo por cidadãos, jornalistas e organizações da sociedade civil.

\section{Comparação com benchmarks acadêmicos}

A equivalência funcional de 72,5\% (Prompt B) e 62,5\% (E1) \textbf{não é diretamente comparável} à EX de Spider ou BIRD, pelas razões seguintes:

\begin{itemize}
    \item \textbf{Métrica distinta}: Spider/BIRD reportam EX estrita (multiconjunto de pares coluna--valor); este trabalho adota EF com mapeamento de colunas de resposta e tolerância $\tau = 0{,}5\%$ em valores fiscais. No E1, EX estrita foi de apenas 10\%, enquanto EF foi de 62,5\% --- evidência de que EX estrita subestimaria respostas úteis ao cidadão;
    \item \textbf{Domínio e escala}: Spider reúne 200 bases em inglês com esquemas sintéticos; BIRD incorpora bases reais volumosas com ruído. O SIM opera sobre esquema SICOM único, em português, com regras contábeis locais --- tarefa mais estreita em diversidade de esquema, porém mais fiel ao uso municipal;
    \item \textbf{Ordem de grandeza}: em BIRD, GPT-4 reporta EX na faixa de 50--55\% \cite{Li2024BIRD}; métodos decomposicionais (DIN-SQL) elevam para $\approx$60\% \cite{Pourreza2024DIN}. A EF de 72,5\% do SIM situa-se na mesma ordem de magnitude, com ressalva de que EF é métrica mais permissiva que EX;
    \item \textbf{Amostra}: o conjunto de referência SQL tem $n=80$ casos correlacionados (10 intenções $\times$ variações), não 80 perguntas independentes --- intervalos de confiança amplos (IC EF do Prompt B: 61,9--81,1\%) refletem incerteza amostral.
\end{itemize}

Em síntese, o SIM não compete em \textit{leaderboard} acadêmico, mas demonstra viabilidade funcional em dados municipais reais com latência compatível a WhatsApp --- objetivo distinto dos \textit{benchmarks} generalistas.

\section{Achados empíricos principais}

Três achados estruturam a discussão, sem repetir numericamente o Capítulo~4:

\begin{enumerate}
    \item \textbf{Engenharia de contexto supera troca de modelo}: Prompt B (+10 pp EF) superou GPT-4o (E4) em EF com fração da latência --- reforçando que regras de domínio SICOM são alavanca principal neste cenário;
    \item \textbf{Recuperação híbrida é necessária para atos numerados}: no estrato \texttt{com\_numero}, busca densa isolada falhou (3,1\% Recall@5); lookup por número/ano no cabeçalho do \textit{chunk} resolve o caso de uso mais frequente em Diário Oficial;
    \item \textbf{Falhas residuais concentram-se em cruzamentos}: EF em consultas ``difíceis'' (57,1\%) e cruzamentos investigativos (30\%) indica que \textit{joins} entre SICOM e CGU por nome ou documento permanecem desafio mesmo com Prompt B.
\end{enumerate}

\section{Configuração adotada na instância Caeté}

Com base nos três lotes experimentais, a instância Caeté adota: (i)~GPT-4.1-mini com Prompt B para Text-to-SQL; (ii)~recuperação híbrida com deduplicação e filtro por tipo de ato para o subsistema RAG. Modelos alternativos e \textit{prompt} estendido foram descartados por desempenho inferior ou latência incompatível com canal conversacional síncrono.

\section{Confiabilidade, risco e segurança}

Mesmo na configuração final (EF 72,5\%), aproximadamente um quarto das consultas do conjunto de referência não produz grandezas funcionalmente equivalentes ao gabarito. Em canal público voltado ao cidadão, isso implica risco de respostas fiscalmente incorretas --- especialmente em cruzamentos investigativos e índices derivados. Mitigações recomendadas incluem: (i)~exibir a consulta SQL gerada ou um resumo auditável quando a pergunta envolve valores monetários; (ii)~mensagens de incerteza quando a confiança do resultado é baixa; e (iii)~encaminhamento a fontes oficiais (portal de transparência, Diário Oficial).

Quanto à segurança, um canal conversacional público está exposto a \textit{prompt injection} e, indiretamente, a SQL gerada sobre base local. O uso de API com permissões restritas, validação sintática antes da execução e ausência de comandos de escrita no DuckDB reduzem, mas não eliminam, o risco. A camada de resposta em linguagem natural pode ainda alucinar interpretações sobre resultados corretos --- reforçando a necessidade de transparência do dado tabular subjacente.

\section{Limitações}
\label{sec:limitacoes}

\begin{itemize}
    \item O conjunto SQL contém 80 casos derivados de 10 intenções analíticas com variações paramétricas correlacionadas;
    \item O ganho do Prompt B foi medido no mesmo conjunto usado para diagnosticar falhas; as regras são portáveis ao padrão SICOM, mas a magnitude do ganho pode variar em outras instâncias (McNemar: $p = 0{,}146$);
    \item Conjunto de referência de uma única instância (Caeté);
    \item Lacuna de folha de pagamento 2020--2022 na base local;
    \item Gabarito manual pode não cobrir todas as SQL semanticamente corretas;
    \item Definição automatizada de colunas de resposta em 75 dos 80 itens SQL;
    \item Julgamento automático (JAR) não discriminante na amostra piloto ($n=30$);
    \item Demonstração WhatsApp sem métricas quantitativas de adoção em produção.
\end{itemize}

\section{Implicações práticas}

O SIM demonstra viabilidade técnica de consulta conversacional a dados municipais abertos, com arquitetura replicável para outras prefeituras que operem sobre o padrão SICOM. Os resultados empíricos subsidiam a escolha de modelo, engenharia de contexto e estratégia de recuperação documental em instâncias municipais com perfil de dados semelhante ao de Caeté.

\chapter{\MakeUppercase{Conclusão}}

Este trabalho apresentou o \textbf{SIM --- Sistema de Informações Municipais}, plataforma conversacional multi-instância para acesso a dados de transparência pública, com \textbf{instância piloto em Caeté (MG)} integrando recuperação documental sobre o Diário Oficial e consulta analítica via Text-to-SQL sobre bases SICOM/CGU.

A arquitetura híbrida --- orquestração por LLM, RAG e DuckDB --- demonstra viabilidade técnica de reduzir a barreira entre o cidadão e dados municipais abertos. A validação com 80 consultas de referência comprovou geração confiável de SQL (VSR de 97,5\%) e equivalência funcional (EF) em 62,5\% dos casos no baseline; o refinamento do \textit{prompt} (v2.1) elevou a EF para 72,5\% com VSR integral. A validação piloto do RAG ($N_{\text{RAG}} = 8$) obteve Recall@5 de 100\% e MRR de 63,33\% na busca densa; melhorias de recuperação (deduplicação, busca híbrida por número de ato e filtro por tipo) elevaram o MRR a 87,5\% no mesmo \textit{golden}.

As contribuições incluem: (i)~conceito SIM multi-instância e arquitetura replicável para GovTech municipal; (ii)~\textit{pipeline} ETL e exposição analítica sobre dados SICOM/CGU; (iii)~protocolo de validação reprodutível versionado em \texttt{avaliacao/}; (iv)~evidências empíricas para escolha de modelo e tamanho de contexto; (v)~validação complementar da recuperação documental; e (vi)~repositório público para replicação.

Como trabalhos futuros, propõe-se: ampliação do \textit{golden} RAG e julgamento de respostas em linguagem natural; revisão humana do conjunto \textit{golden} SQL; testes com usuários reais; replicação do SIM em outros municípios; e autocorreção iterativa de consultas SQL.


%==============================================================================
% Bibliografia

% Carrega o estilo para referências ABNT
\bibliographystyle{abntex2-alf}            

% Alterando informações do estilo de citação.
% Como estamos usando uma versão alterada, essas opções são necessárias.
% Et al. em itálico
\citeoption{abnt-etal-text=it}

% Até 3 autores antes do et al.
\citeoption{abnt-etal-cite=3}

% Colocar s.d. na falta de um ano de publicação
\citeoption{abnt-missing-year=sd}

% Altera o estilo de citação para (Autor, ano)
% A opção abnt-cite-style=AuthorYEAR está quebrada; essa é equivalente.
\citeoption{abnt-nbr10520=1988}
\citeoption{abnt-and-type=e}


% Comando necessário para corrigir o título das referências.
% Para o estilo ABNT, esse comando SEMPRE deve ser
% chamado antes do \bibliography
\referencias
% Adiciona as referência em bibtex 
\bibliography{refbib}

%==============================================================================
% Outros elementos pós-textuais

% Glossário
\glossario

% Anexos e apêndices
% Os anexos devem ser includios antes
% e utilizar \input ao invés de \include, para não criar uma nova página
% Anexos
%\input{anexos/anexo1}
%\input{anexos/anexo2}
\apendice{Estrutura e Especificação dos Artefatos Reprodutíveis}\label{ap:artefatos}

Este apêndice detalha a organização dos artefatos computacionais disponibilizados no repositório público do projeto. O módulo de avaliação (\texttt{avaliacao/}) é estruturado em \textit{runners} de experimento, conjuntos de referência versionados, variantes de \textit{prompt} e saídas por experimento --- elementos omitidos do corpo textual para preservar o registro acadêmico, mas necessários à replicação dos resultados reportados nos Capítulos~3 e~4.

\section{Correspondência nomenclatura e arquivos}

\begin{table}[htbp]
\centering
\caption{Mapeamento entre nomes usados no corpo e arquivos versionados}
\label{tab:map_artefatos}
\small
\begin{tabular}{p{6.5cm}p{6.5cm}}
\toprule
\textbf{Nome no corpo do texto} & \textbf{Arquivo versionado} \\
\midrule
Conjunto de referência SQL (80 casos) & \texttt{golden\_dataset\_v1.0.csv} \\
Conjunto de referência RAG ampliado ($n=40$) & \texttt{golden\_rag\_v1.3.csv} \\
Amostra temática curada ($n=8$) & \texttt{golden\_rag\_v1.1.csv} \\
\bottomrule
\end{tabular}
\end{table}

\section{Estrutura principal}

\begin{itemize}
    \item \texttt{avaliacao/} --- módulo de avaliação reprodutível (SQL e RAG);
    \item \texttt{avaliacao/eval/} --- \textit{runners} de experimento e biblioteca de métricas;
    \item \texttt{avaliacao/dados/golden/} --- conjunto de referência SQL (80 casos);
    \item \texttt{avaliacao/dados/golden\_rag/} --- conjunto de referência RAG ($n=40$);
    \item \texttt{avaliacao/experimentos/} --- saídas versionadas por experimento (\texttt{checkpoint.jsonl}, \texttt{metrics\_summary\_v3.json});
    \item \texttt{avaliacao/eval/prompts/} --- variantes de \textit{prompt} (compacto, estendido, regras de domínio);
    \item \texttt{scripts\_python/} --- ETL e API analítica da instância piloto;
    \item \texttt{docs/} --- documentação técnica, métricas e textos acadêmicos.
\end{itemize}

\section{Metadados do conjunto SQL}

\begin{itemize}
    \item Arquivo principal: \texttt{golden\_dataset\_v1.0.csv} (conjunto de referência SQL, 80 casos);
    \item Campo \texttt{colunas\_resposta}: grandezas avaliadas na equivalência funcional;
    \item Campo \texttt{origem}: \texttt{base\_10} (intenção analítica) ou \texttt{derivacao} (variação paramétrica);
    \item Resultados de referência: \texttt{reference\_results/NNN.json}.
\end{itemize}

\section{Metadados do conjunto RAG}

\begin{itemize}
    \item Arquivo principal: \texttt{golden\_rag\_v1.3.csv} (conjunto de referência RAG ampliado, $n=40$);
    \item Campo \texttt{doc\_id\_esperado}: rótulo canônico do ato no Diário Oficial;
    \item Campo \texttt{estrato}: \texttt{com\_numero} ou \texttt{tematico};
    \item Índice Qdrant: \textit{collection} \texttt{jornais\_caete}; correspondência por \texttt{page\_content}.
\end{itemize}

\section{Execução da avaliação}

A reprodução dos experimentos é automatizada pelo script \texttt{run\_metrics\_pipeline.ps1}, no diretório \texttt{avaliacao/}. Na forma padrão, o \textit{pipeline} executa a regressão SQL (Lote~I, experimento E1); com o parâmetro \texttt{-WithRag}, inclui também a avaliação do subsistema de recuperação documental (Lote~III). A configuração central reside em \texttt{avaliacao/eval/config.yaml}. Os requisitos de ambiente são: instância DuckDB local sobre as bases consolidadas, serviço Qdrant acessível na porta \texttt{6333} e chave de API OpenAI para \textit{embeddings} na avaliação RAG. O script completo, dependências e instruções de instalação estão versionados no repositório público do projeto.


% Índice
\indice


%==============================================================================
% Fim do documento
%==============================================================================
\end{document}
